
\documentclass{article}
\usepackage{colortbl}
\usepackage{makecell}
\usepackage{multirow}
\usepackage{supertabular}

\begin{document}

\newcounter{utterance}

\centering \large Interaction Transcript for game `dond', experiment `coop\_en', episode 16 with Qwen3{-}32B{-}t0.0.
\vspace{24pt}

{ \footnotesize  \setcounter{utterance}{1}
\setlength{\tabcolsep}{0pt}
\begin{supertabular}{c@{$\;$}|p{.15\linewidth}@{}p{.15\linewidth}p{.15\linewidth}p{.15\linewidth}p{.15\linewidth}p{.15\linewidth}}
    \# & $\;$A & \multicolumn{4}{c}{Game Master} & $\;\:$B\\
    \hline

    \theutterance \stepcounter{utterance}  
    & & \multicolumn{4}{p{0.6\linewidth}}{
        \cellcolor[rgb]{0.9,0.9,0.9}{
            \makecell[{{p{\linewidth}}}]{
                \texttt{\tiny{[P1$\langle$GM]}}
                \texttt{You are playing a cooperative negotiation game in which you have to agree on how to divide a set of items among you and another player.} \\
\\ 
\texttt{Rules:} \\
\texttt{(a) You and the other player are given a set of items. Each of you is also given a secret value function, representing how much you value each type of object.} \\
\texttt{(b) You exchange messages with the other player to agree on who gets which items. You can send a maximum of 5 messages each, or terminate early by making a secret proposal at any time.} \\
\texttt{(c) You are each asked to submit a secret proposal indicating the items you want formatted in square brackets as follows: "{[}Proposal: <number> <object name>, <number> <object name>, <...>{]}"} \\
\texttt{(d) If your proposals are complementary, i.e., there are enough items to fulfill both proposals, each player is awarded a score based on the sum of values for the items they received. Otherwise, both of you get zero points.} \\
\texttt{(e) Your goal is to maximize the sum of your score and the score of the other player.} \\
\\ 
\texttt{Let us start.} \\
\\ 
\texttt{The set of available items is:} \\
\\ 
\texttt{1 bucket, 5 books, 1 mirror} \\
\\ 
\texttt{Your secret value function is:} \\
\\ 
\texttt{bucket: 4, book: 1, mirror: 1} \\
\\ 
\texttt{IMPORTANT: Your messages, unless it is the secret proposal, are directly transmitted to the other player, so do not include any response to the rules or text announcing your message. To make a secret proposal, use the indicated format. Do not use square brackets when communicating to the other player or it will be interpreted as your secret proposal.} \\
            }
        }
    }
    & \\ \\

    \theutterance \stepcounter{utterance}  
    & \multicolumn{4}{p{0.6\linewidth}}{
        \cellcolor[rgb]{0.9,0.9,0.9}{
            \makecell[{{p{\linewidth}}}]{
                \texttt{\tiny{[P1$\rangle$GM]}}
                \texttt{<think>} \\
\texttt{Okay, let's see. I need to divide the items between me and another player. The items are 1 bucket, 5 books, and 1 mirror. My value function is bucket:4, book:1, mirror:1. The other player's value function is secret, so I don't know what they value. The goal is to maximize the total score for both of us.} \\
\\ 
\texttt{First, I should figure out what the other player might value. Since I don't have their value function, I need to make some assumptions. Maybe they value the bucket more, or maybe the mirror? Or perhaps the books? Since there are 5 books, maybe they have a higher value for books. But without knowing, I need to find a way to split the items so that both of us can get what we value most without overlapping too much.} \\
\\ 
\texttt{My own highest value is the bucket at 4. The mirror and books are each 1. So I should probably want the bucket for myself. But if the other player also values the bucket, then we might both want it, leading to a conflict. However, since there's only one bucket, we can't both have it. So maybe I should propose taking the bucket and some books, and let the other player take the mirror and the remaining books. But how many books?} \\
\\ 
\texttt{Let me think. If I take the bucket (4 points) and, say, 2 books (2 points), that's 6 points. The other player would get 3 books (3 points) and the mirror. If their mirror is worth more to them, maybe that's a good split. Alternatively, if they value the mirror less, maybe they'd prefer more books. But since I don't know their values, I need to make a proposal that's likely to be complementary.} \\
\\ 
\texttt{Alternatively, maybe I should take the bucket and all the books. But that would leave the other player with only the mirror. If their mirror is worth a lot to them, maybe that's okay. But if they value books more, then they might not accept. But since the proposal has to be complementary, if I take the bucket and all 5 books, the other player can only take the mirror. So that would be complementary. But my total would be 4 + 5*1 = 9, and their mirror's value depends on their function. If their mirror is, say, 5, then total would} \\
            }
        }
    }
    & & \\ \\

    \theutterance \stepcounter{utterance}  
    & & & \multicolumn{2}{p{0.3\linewidth}}{
        \cellcolor[rgb]{0.9,0.9,0.9}{
            \makecell[{{p{\linewidth}}}]{
                \texttt{\tiny{[GM$|$GM]}}
                \texttt{continue} \\
            }
        }
    }
    & & \\ \\

    \theutterance \stepcounter{utterance}  
    & & & \multicolumn{2}{p{0.3\linewidth}}{
        \cellcolor[rgb]{0.9,0.9,0.9}{
            \makecell[{{p{\linewidth}}}]{
                \texttt{\tiny{[GM$|$GM]}}
                \texttt{<think>} \\
\texttt{Okay, let's see. I need to divide the items between me and another player. The items are 1 bucket, 5 books, and 1 mirror. My value function is bucket:4, book:1, mirror:1. The other player's value function is secret, so I don't know what they value. The goal is to maximize the total score for both of us.} \\
\\ 
\texttt{First, I should figure out what the other player might value. Since I don't have their value function, I need to make some assumptions. Maybe they value the bucket more, or maybe the mirror? Or perhaps the books? Since there are 5 books, maybe they have a higher value for books. But without knowing, I need to find a way to split the items so that both of us can get what we value most without overlapping too much.} \\
\\ 
\texttt{My own highest value is the bucket at 4. The mirror and books are each 1. So I should probably want the bucket for myself. But if the other player also values the bucket, then we might both want it, leading to a conflict. However, since there's only one bucket, we can't both have it. So maybe I should propose taking the bucket and some books, and let the other player take the mirror and the remaining books. But how many books?} \\
\\ 
\texttt{Let me think. If I take the bucket (4 points) and, say, 2 books (2 points), that's 6 points. The other player would get 3 books (3 points) and the mirror. If their mirror is worth more to them, maybe that's a good split. Alternatively, if they value the mirror less, maybe they'd prefer more books. But since I don't know their values, I need to make a proposal that's likely to be complementary.} \\
\\ 
\texttt{Alternatively, maybe I should take the bucket and all the books. But that would leave the other player with only the mirror. If their mirror is worth a lot to them, maybe that's okay. But if they value books more, then they might not accept. But since the proposal has to be complementary, if I take the bucket and all 5 books, the other player can only take the mirror. So that would be complementary. But my total would be 4 + 5*1 = 9, and their mirror's value depends on their function. If their mirror is, say, 5, then total would} \\
            }
        }
    }
    & & \\ \\

    \theutterance \stepcounter{utterance}  
    & & \multicolumn{4}{p{0.6\linewidth}}{
        \cellcolor[rgb]{0.9,0.9,0.9}{
            \makecell[{{p{\linewidth}}}]{
                \texttt{\tiny{[GM$\rangle$P2]}}
                \texttt{You are playing a cooperative negotiation game in which you have to agree on how to divide a set of items among you and another player.} \\
\\ 
\texttt{Rules:} \\
\texttt{(a) You and the other player are given a set of items. Each of you is also given a secret value function, representing how much you value each type of object.} \\
\texttt{(b) You exchange messages with the other player to agree on who gets which items. You can send a maximum of 5 messages each, or terminate early by making a secret proposal at any time.} \\
\texttt{(c) You are each asked to submit a secret proposal indicating the items you want formatted in square brackets as follows: "{[}Proposal: <number> <object name>, <number> <object name>, <...>{]}"} \\
\texttt{(d) If your proposals are complementary, i.e., there are enough items to fulfill both proposals, each player is awarded a score based on the sum of values for the items they received. Otherwise, both of you get zero points.} \\
\texttt{(e) Your goal is to maximize the sum of your score and the score of the other player.} \\
\\ 
\texttt{Let us start.} \\
\\ 
\texttt{The set of available items is:} \\
\\ 
\texttt{1 bucket, 5 books, 1 mirror} \\
\\ 
\texttt{Your secret value function is:} \\
\\ 
\texttt{bucket: 2, book: 0, mirror: 8} \\
\\ 
\texttt{IMPORTANT: Your messages, unless it is the secret proposal, are directly transmitted to the other player, so do not include any response to the rules or text announcing your message. To make a secret proposal, use the indicated format. Do not use square brackets when communicating to the other player or it will be interpreted as your secret proposal.} \\
\\ 
\\ 
\texttt{<think>} \\
\texttt{Okay, let's see. I need to divide the items between me and another player. The items are 1 bucket, 5 books, and 1 mirror. My value function is bucket:4, book:1, mirror:1. The other player's value function is secret, so I don't know what they value. The goal is to maximize the total score for both of us.} \\
\\ 
\texttt{First, I should figure out what the other player might value. Since I don't have their value function, I need to make some assumptions. Maybe they value the bucket more, or maybe the mirror? Or perhaps the books? Since there are 5 books, maybe they have a higher value for books. But without knowing, I need to find a way to split the items so that both of us can get what we value most without overlapping too much.} \\
\\ 
\texttt{My own highest value is the bucket at 4. The mirror and books are each 1. So I should probably want the bucket for myself. But if the other player also values the bucket, then we might both want it, leading to a conflict. However, since there's only one bucket, we can't both have it. So maybe I should propose taking the bucket and some books, and let the other player take the mirror and the remaining books. But how many books?} \\
\\ 
\texttt{Let me think. If I take the bucket (4 points) and, say, 2 books (2 points), that's 6 points. The other player would get 3 books (3 points) and the mirror. If their mirror is worth more to them, maybe that's a good split. Alternatively, if they value the mirror less, maybe they'd prefer more books. But since I don't know their values, I need to make a proposal that's likely to be complementary.} \\
\\ 
\texttt{Alternatively, maybe I should take the bucket and all the books. But that would leave the other player with only the mirror. If their mirror is worth a lot to them, maybe that's okay. But if they value books more, then they might not accept. But since the proposal has to be complementary, if I take the bucket and all 5 books, the other player can only take the mirror. So that would be complementary. But my total would be 4 + 5*1 = 9, and their mirror's value depends on their function. If their mirror is, say, 5, then total would} \\
            }
        }
    }
    & \\ \\

    \theutterance \stepcounter{utterance}  
    & & & \multicolumn{4}{p{0.6\linewidth}}{
        \cellcolor[rgb]{0.9,0.9,0.9}{
            \makecell[{{p{\linewidth}}}]{
                \texttt{\tiny{[GM$\langle$P2]}}
                \texttt{<think>} \\
\texttt{Okay, let's tackle this problem step by step. So, the goal is to divide the items between me and another player in a way that maximizes the sum of both our scores. The items available are 1 bucket, 5 books, and 1 mirror. My value function is bucket: 4, book: 1, mirror: 1. The other player's values are secret, so I need to make an educated guess or find a way to split the items that likely accommodates their preferences without overlap.} \\
\\ 
\texttt{First, I need to consider what the other player might value. Since I don't have their value function, I have to think about possible scenarios. The other player could value the bucket, the mirror, or the books more. Let's explore each possibility.} \\
\\ 
\texttt{If the other player values the bucket highly, then I should probably let them take the bucket, and I can take the mirror and some books. But my own value for the bucket is 4, which is higher than the mirror and books. However, if the other player's value for the bucket is even higher, maybe it's better for both of us if they take it. But since I don't know, this is a risk.} \\
\\ 
\texttt{Alternatively, if the other player values the mirror more, then I could take the bucket and some books, and they take the mirror and the remaining books. That way, I get my high{-}value item (bucket) and some books, and they get the mirror and more books. Since there are 5 books, splitting them between us might be necessary. For example, I take the bucket and 2 books, they take the mirror and 3 books. My total would be 4 + 2*1 = 6, and their total would depend on their values for mirror and books. If their mirror is worth, say, 5, and books 1 each, then they get 5 + 3 = 8, total combined 14. If I took the bucket and all books, my total would be 4 +5=9, and they get mirror. If their mirror is 5, total is 14 as well. So maybe similar.} \\
\\ 
\texttt{But if the other player values books more, maybe they want as many as possible. If I take the bucket and 1 book, they can take 4 books and the mirror. My score would be 4 +1=5, theirs would be} \\
            }
        }
    }
     \\ \\

    \theutterance \stepcounter{utterance}  
    & & & \multicolumn{2}{p{0.3\linewidth}}{
        \cellcolor[rgb]{0.9,0.9,0.9}{
            \makecell[{{p{\linewidth}}}]{
                \texttt{\tiny{[GM$|$GM]}}
                \texttt{continue} \\
            }
        }
    }
    & & \\ \\

    \theutterance \stepcounter{utterance}  
    & & & \multicolumn{2}{p{0.3\linewidth}}{
        \cellcolor[rgb]{0.9,0.9,0.9}{
            \makecell[{{p{\linewidth}}}]{
                \texttt{\tiny{[GM$|$GM]}}
                \texttt{<think>} \\
\texttt{Okay, let's tackle this problem step by step. So, the goal is to divide the items between me and another player in a way that maximizes the sum of both our scores. The items available are 1 bucket, 5 books, and 1 mirror. My value function is bucket: 4, book: 1, mirror: 1. The other player's values are secret, so I need to make an educated guess or find a way to split the items that likely accommodates their preferences without overlap.} \\
\\ 
\texttt{First, I need to consider what the other player might value. Since I don't have their value function, I have to think about possible scenarios. The other player could value the bucket, the mirror, or the books more. Let's explore each possibility.} \\
\\ 
\texttt{If the other player values the bucket highly, then I should probably let them take the bucket, and I can take the mirror and some books. But my own value for the bucket is 4, which is higher than the mirror and books. However, if the other player's value for the bucket is even higher, maybe it's better for both of us if they take it. But since I don't know, this is a risk.} \\
\\ 
\texttt{Alternatively, if the other player values the mirror more, then I could take the bucket and some books, and they take the mirror and the remaining books. That way, I get my high{-}value item (bucket) and some books, and they get the mirror and more books. Since there are 5 books, splitting them between us might be necessary. For example, I take the bucket and 2 books, they take the mirror and 3 books. My total would be 4 + 2*1 = 6, and their total would depend on their values for mirror and books. If their mirror is worth, say, 5, and books 1 each, then they get 5 + 3 = 8, total combined 14. If I took the bucket and all books, my total would be 4 +5=9, and they get mirror. If their mirror is 5, total is 14 as well. So maybe similar.} \\
\\ 
\texttt{But if the other player values books more, maybe they want as many as possible. If I take the bucket and 1 book, they can take 4 books and the mirror. My score would be 4 +1=5, theirs would be} \\
            }
        }
    }
    & & \\ \\

    \theutterance \stepcounter{utterance}  
    & & \multicolumn{4}{p{0.6\linewidth}}{
        \cellcolor[rgb]{0.9,0.9,0.9}{
            \makecell[{{p{\linewidth}}}]{
                \texttt{\tiny{[P1$\langle$GM]}}
                \texttt{<think>} \\
\texttt{Okay, let's tackle this problem step by step. So, the goal is to divide the items between me and another player in a way that maximizes the sum of both our scores. The items available are 1 bucket, 5 books, and 1 mirror. My value function is bucket: 4, book: 1, mirror: 1. The other player's values are secret, so I need to make an educated guess or find a way to split the items that likely accommodates their preferences without overlap.} \\
\\ 
\texttt{First, I need to consider what the other player might value. Since I don't have their value function, I have to think about possible scenarios. The other player could value the bucket, the mirror, or the books more. Let's explore each possibility.} \\
\\ 
\texttt{If the other player values the bucket highly, then I should probably let them take the bucket, and I can take the mirror and some books. But my own value for the bucket is 4, which is higher than the mirror and books. However, if the other player's value for the bucket is even higher, maybe it's better for both of us if they take it. But since I don't know, this is a risk.} \\
\\ 
\texttt{Alternatively, if the other player values the mirror more, then I could take the bucket and some books, and they take the mirror and the remaining books. That way, I get my high{-}value item (bucket) and some books, and they get the mirror and more books. Since there are 5 books, splitting them between us might be necessary. For example, I take the bucket and 2 books, they take the mirror and 3 books. My total would be 4 + 2*1 = 6, and their total would depend on their values for mirror and books. If their mirror is worth, say, 5, and books 1 each, then they get 5 + 3 = 8, total combined 14. If I took the bucket and all books, my total would be 4 +5=9, and they get mirror. If their mirror is 5, total is 14 as well. So maybe similar.} \\
\\ 
\texttt{But if the other player values books more, maybe they want as many as possible. If I take the bucket and 1 book, they can take 4 books and the mirror. My score would be 4 +1=5, theirs would be} \\
            }
        }
    }
    & \\ \\

    \theutterance \stepcounter{utterance}  
    & \multicolumn{4}{p{0.6\linewidth}}{
        \cellcolor[rgb]{0.9,0.9,0.9}{
            \makecell[{{p{\linewidth}}}]{
                \texttt{\tiny{[P1$\rangle$GM]}}
                \texttt{<think>} \\
\texttt{Okay, let's see. I need to figure out how to split the items so that both me and the other player get what we value most without overlapping. The items are 1 bucket, 5 books, and 1 mirror. My values are bucket:4, book:1, mirror:1. The other player's values are secret, so I have to make an educated guess.} \\
\\ 
\texttt{First, I should consider that the other player might value the bucket, mirror, or books more. Since there's only one bucket and one mirror, but five books, maybe the other player wants the bucket or mirror. But I don't know. My highest value is the bucket, so I want that. But if the other player also wants the bucket, we can't both have it. So maybe I should propose taking the bucket and some books, and let them take the mirror and the rest of the books. That way, we don't conflict on the bucket and mirror. Let's say I take the bucket and 2 books (4 + 2 = 6 points), and they take the mirror and 3 books. If their mirror is worth more to them than the books, that's good. If not, maybe they'd prefer more books. But since I don't know their values, this split might be safe. Alternatively, if I take the bucket and all 5 books, I get 4 +5=9, and they get the mirror. If their mirror is worth a lot, that's okay. But if they value books more, they might not get much. But since the proposals need to be complementary, if I take the bucket and all books, they can only take the mirror. So that's a valid split. But maybe the other player would prefer to have some books. Hmm. Let me think. If I take the bucket and 3 books, I get 4 +3=7, and they get 2 books and the mirror. If their mirror is worth, say, 3, and books 1 each, they get 3 +2=5. Total combined 12. If I take bucket and 2 books (6) and they take mirror and 3 books (3+3=6), total 12. Same. If I take bucket and 1 book (5) and they take mirror and 4 books (3+4=7), total 12. So} \\
            }
        }
    }
    & & \\ \\

    \theutterance \stepcounter{utterance}  
    & & & \multicolumn{2}{p{0.3\linewidth}}{
        \cellcolor[rgb]{0.9,0.9,0.9}{
            \makecell[{{p{\linewidth}}}]{
                \texttt{\tiny{[GM$|$GM]}}
                \texttt{continue} \\
            }
        }
    }
    & & \\ \\

    \theutterance \stepcounter{utterance}  
    & & & \multicolumn{2}{p{0.3\linewidth}}{
        \cellcolor[rgb]{0.9,0.9,0.9}{
            \makecell[{{p{\linewidth}}}]{
                \texttt{\tiny{[GM$|$GM]}}
                \texttt{<think>} \\
\texttt{Okay, let's see. I need to figure out how to split the items so that both me and the other player get what we value most without overlapping. The items are 1 bucket, 5 books, and 1 mirror. My values are bucket:4, book:1, mirror:1. The other player's values are secret, so I have to make an educated guess.} \\
\\ 
\texttt{First, I should consider that the other player might value the bucket, mirror, or books more. Since there's only one bucket and one mirror, but five books, maybe the other player wants the bucket or mirror. But I don't know. My highest value is the bucket, so I want that. But if the other player also wants the bucket, we can't both have it. So maybe I should propose taking the bucket and some books, and let them take the mirror and the rest of the books. That way, we don't conflict on the bucket and mirror. Let's say I take the bucket and 2 books (4 + 2 = 6 points), and they take the mirror and 3 books. If their mirror is worth more to them than the books, that's good. If not, maybe they'd prefer more books. But since I don't know their values, this split might be safe. Alternatively, if I take the bucket and all 5 books, I get 4 +5=9, and they get the mirror. If their mirror is worth a lot, that's okay. But if they value books more, they might not get much. But since the proposals need to be complementary, if I take the bucket and all books, they can only take the mirror. So that's a valid split. But maybe the other player would prefer to have some books. Hmm. Let me think. If I take the bucket and 3 books, I get 4 +3=7, and they get 2 books and the mirror. If their mirror is worth, say, 3, and books 1 each, they get 3 +2=5. Total combined 12. If I take bucket and 2 books (6) and they take mirror and 3 books (3+3=6), total 12. Same. If I take bucket and 1 book (5) and they take mirror and 4 books (3+4=7), total 12. So} \\
            }
        }
    }
    & & \\ \\

    \theutterance \stepcounter{utterance}  
    & & \multicolumn{4}{p{0.6\linewidth}}{
        \cellcolor[rgb]{0.9,0.9,0.9}{
            \makecell[{{p{\linewidth}}}]{
                \texttt{\tiny{[GM$\rangle$P2]}}
                \texttt{<think>} \\
\texttt{Okay, let's see. I need to figure out how to split the items so that both me and the other player get what we value most without overlapping. The items are 1 bucket, 5 books, and 1 mirror. My values are bucket:4, book:1, mirror:1. The other player's values are secret, so I have to make an educated guess.} \\
\\ 
\texttt{First, I should consider that the other player might value the bucket, mirror, or books more. Since there's only one bucket and one mirror, but five books, maybe the other player wants the bucket or mirror. But I don't know. My highest value is the bucket, so I want that. But if the other player also wants the bucket, we can't both have it. So maybe I should propose taking the bucket and some books, and let them take the mirror and the rest of the books. That way, we don't conflict on the bucket and mirror. Let's say I take the bucket and 2 books (4 + 2 = 6 points), and they take the mirror and 3 books. If their mirror is worth more to them than the books, that's good. If not, maybe they'd prefer more books. But since I don't know their values, this split might be safe. Alternatively, if I take the bucket and all 5 books, I get 4 +5=9, and they get the mirror. If their mirror is worth a lot, that's okay. But if they value books more, they might not get much. But since the proposals need to be complementary, if I take the bucket and all books, they can only take the mirror. So that's a valid split. But maybe the other player would prefer to have some books. Hmm. Let me think. If I take the bucket and 3 books, I get 4 +3=7, and they get 2 books and the mirror. If their mirror is worth, say, 3, and books 1 each, they get 3 +2=5. Total combined 12. If I take bucket and 2 books (6) and they take mirror and 3 books (3+3=6), total 12. Same. If I take bucket and 1 book (5) and they take mirror and 4 books (3+4=7), total 12. So} \\
            }
        }
    }
    & \\ \\

    \theutterance \stepcounter{utterance}  
    & & & \multicolumn{4}{p{0.6\linewidth}}{
        \cellcolor[rgb]{0.9,0.9,0.9}{
            \makecell[{{p{\linewidth}}}]{
                \texttt{\tiny{[GM$\langle$P2]}}
                \texttt{<think>} \\
\texttt{Okay, let's try to figure out the best way to split the items. The available items are 1 bucket, 5 books, and 1 mirror. My value function is bucket:4, book:1, mirror:1. The other player's values are secret, so I need to make a proposal that likely maximizes the total score without knowing their preferences.} \\
\\ 
\texttt{First, I should consider that the other player might value the bucket, mirror, or books more. Since there's only one bucket and one mirror, but five books, it's possible they might want the bucket or mirror. My highest value is the bucket, so I want that. But if the other player also wants the bucket, we can't both have it. So maybe I should propose taking the bucket and some books, and let them take the mirror and the rest of the books. That way, we don't conflict on the bucket and mirror. Let's say I take the bucket and 2 books (4 + 2 = 6 points), and they take the mirror and 3 books. If their mirror is worth more to them than the books, that's good. If not, maybe they'd prefer more books. But since I don't know their values, this split might be safe. Alternatively, if I take the bucket and all 5 books, I get 4 +5=9, and they get the mirror. If their mirror is worth a lot, that's okay. But if they value books more, they might not get much. But since the proposals need to be complementary, if I take the bucket and all books, they can only take the mirror. So that's a valid split. But maybe the other player would prefer to have some books. Hmm. Let me think. If I take the bucket and 3 books, I get 4 +3=7, and they get 2 books and the mirror. If their mirror is worth, say, 3, and books 1 each, they get 3 +2=5. Total combined 12. If I take bucket and 2 books (6) and they take mirror and 3 books (3+3=6), total 12. Same. If I take bucket and 1 book (5) and they take mirror and 4 books (3+4=7), total 12. So regardless of how I split the books, the} \\
            }
        }
    }
     \\ \\

    \theutterance \stepcounter{utterance}  
    & & & \multicolumn{2}{p{0.3\linewidth}}{
        \cellcolor[rgb]{0.9,0.9,0.9}{
            \makecell[{{p{\linewidth}}}]{
                \texttt{\tiny{[GM$|$GM]}}
                \texttt{continue} \\
            }
        }
    }
    & & \\ \\

    \theutterance \stepcounter{utterance}  
    & & & \multicolumn{2}{p{0.3\linewidth}}{
        \cellcolor[rgb]{0.9,0.9,0.9}{
            \makecell[{{p{\linewidth}}}]{
                \texttt{\tiny{[GM$|$GM]}}
                \texttt{<think>} \\
\texttt{Okay, let's try to figure out the best way to split the items. The available items are 1 bucket, 5 books, and 1 mirror. My value function is bucket:4, book:1, mirror:1. The other player's values are secret, so I need to make a proposal that likely maximizes the total score without knowing their preferences.} \\
\\ 
\texttt{First, I should consider that the other player might value the bucket, mirror, or books more. Since there's only one bucket and one mirror, but five books, it's possible they might want the bucket or mirror. My highest value is the bucket, so I want that. But if the other player also wants the bucket, we can't both have it. So maybe I should propose taking the bucket and some books, and let them take the mirror and the rest of the books. That way, we don't conflict on the bucket and mirror. Let's say I take the bucket and 2 books (4 + 2 = 6 points), and they take the mirror and 3 books. If their mirror is worth more to them than the books, that's good. If not, maybe they'd prefer more books. But since I don't know their values, this split might be safe. Alternatively, if I take the bucket and all 5 books, I get 4 +5=9, and they get the mirror. If their mirror is worth a lot, that's okay. But if they value books more, they might not get much. But since the proposals need to be complementary, if I take the bucket and all books, they can only take the mirror. So that's a valid split. But maybe the other player would prefer to have some books. Hmm. Let me think. If I take the bucket and 3 books, I get 4 +3=7, and they get 2 books and the mirror. If their mirror is worth, say, 3, and books 1 each, they get 3 +2=5. Total combined 12. If I take bucket and 2 books (6) and they take mirror and 3 books (3+3=6), total 12. Same. If I take bucket and 1 book (5) and they take mirror and 4 books (3+4=7), total 12. So regardless of how I split the books, the} \\
            }
        }
    }
    & & \\ \\

    \theutterance \stepcounter{utterance}  
    & & \multicolumn{4}{p{0.6\linewidth}}{
        \cellcolor[rgb]{0.9,0.9,0.9}{
            \makecell[{{p{\linewidth}}}]{
                \texttt{\tiny{[P1$\langle$GM]}}
                \texttt{<think>} \\
\texttt{Okay, let's try to figure out the best way to split the items. The available items are 1 bucket, 5 books, and 1 mirror. My value function is bucket:4, book:1, mirror:1. The other player's values are secret, so I need to make a proposal that likely maximizes the total score without knowing their preferences.} \\
\\ 
\texttt{First, I should consider that the other player might value the bucket, mirror, or books more. Since there's only one bucket and one mirror, but five books, it's possible they might want the bucket or mirror. My highest value is the bucket, so I want that. But if the other player also wants the bucket, we can't both have it. So maybe I should propose taking the bucket and some books, and let them take the mirror and the rest of the books. That way, we don't conflict on the bucket and mirror. Let's say I take the bucket and 2 books (4 + 2 = 6 points), and they take the mirror and 3 books. If their mirror is worth more to them than the books, that's good. If not, maybe they'd prefer more books. But since I don't know their values, this split might be safe. Alternatively, if I take the bucket and all 5 books, I get 4 +5=9, and they get the mirror. If their mirror is worth a lot, that's okay. But if they value books more, they might not get much. But since the proposals need to be complementary, if I take the bucket and all books, they can only take the mirror. So that's a valid split. But maybe the other player would prefer to have some books. Hmm. Let me think. If I take the bucket and 3 books, I get 4 +3=7, and they get 2 books and the mirror. If their mirror is worth, say, 3, and books 1 each, they get 3 +2=5. Total combined 12. If I take bucket and 2 books (6) and they take mirror and 3 books (3+3=6), total 12. Same. If I take bucket and 1 book (5) and they take mirror and 4 books (3+4=7), total 12. So regardless of how I split the books, the} \\
            }
        }
    }
    & \\ \\

    \theutterance \stepcounter{utterance}  
    & \multicolumn{4}{p{0.6\linewidth}}{
        \cellcolor[rgb]{0.9,0.9,0.9}{
            \makecell[{{p{\linewidth}}}]{
                \texttt{\tiny{[P1$\rangle$GM]}}
                \texttt{<think>} \\
\texttt{Okay, let's try to figure out the best way to split the items. The available items are 1 bucket, 5 books, and 1 mirror. My value function is bucket:4, book:1, mirror:1. The other player's values are secret, so I need to make a proposal that likely maximizes the total score without knowing their preferences.} \\
\\ 
\texttt{First, I should consider that the other player might value the bucket, mirror, or books more. Since there's only one bucket and one mirror, but five books, it's possible they might want the bucket or mirror. My highest value is the bucket, so I want that. But if the other player also wants the bucket, we can't both have it. So maybe I should propose taking the bucket and some books, and let them take the mirror and the rest of the books. That way, we don't conflict on the bucket and mirror. Let's say I take the bucket and 2 books (4 + 2 = 6 points), and they take the mirror and 3 books. If their mirror is worth more to them than the books, that's good. If not, maybe they'd prefer more books. But since I don't know their values, this split might be safe. Alternatively, if I take the bucket and all 5 books, I get 4 +5=9, and they get the mirror. If their mirror is worth a lot, that's okay. But if they value books more, they might not get much. But since the proposals need to be complementary, if I take the bucket and all books, they can only take the mirror. So that's a valid split. But maybe the other player would prefer to have some books. Hmm. Let me think. If I take the bucket and 3 books, I get 4 +3=7, and they get 2 books and the mirror. If their mirror is worth, say, 3, and books 1 each, they get 3 +2=5. Total combined 12. If I take bucket and 2 books (6) and they take mirror and 3 books (3+3=6), total 12. Same. If I take bucket and 1 book (5) and they take mirror and 4 books (3+4=7), total 12. So regardless of how I split the books, the} \\
            }
        }
    }
    & & \\ \\

    \theutterance \stepcounter{utterance}  
    & & & \multicolumn{2}{p{0.3\linewidth}}{
        \cellcolor[rgb]{0.9,0.9,0.9}{
            \makecell[{{p{\linewidth}}}]{
                \texttt{\tiny{[GM$|$GM]}}
                \texttt{continue} \\
            }
        }
    }
    & & \\ \\

    \theutterance \stepcounter{utterance}  
    & & & \multicolumn{2}{p{0.3\linewidth}}{
        \cellcolor[rgb]{0.9,0.9,0.9}{
            \makecell[{{p{\linewidth}}}]{
                \texttt{\tiny{[GM$|$GM]}}
                \texttt{<think>} \\
\texttt{Okay, let's try to figure out the best way to split the items. The available items are 1 bucket, 5 books, and 1 mirror. My value function is bucket:4, book:1, mirror:1. The other player's values are secret, so I need to make a proposal that likely maximizes the total score without knowing their preferences.} \\
\\ 
\texttt{First, I should consider that the other player might value the bucket, mirror, or books more. Since there's only one bucket and one mirror, but five books, it's possible they might want the bucket or mirror. My highest value is the bucket, so I want that. But if the other player also wants the bucket, we can't both have it. So maybe I should propose taking the bucket and some books, and let them take the mirror and the rest of the books. That way, we don't conflict on the bucket and mirror. Let's say I take the bucket and 2 books (4 + 2 = 6 points), and they take the mirror and 3 books. If their mirror is worth more to them than the books, that's good. If not, maybe they'd prefer more books. But since I don't know their values, this split might be safe. Alternatively, if I take the bucket and all 5 books, I get 4 +5=9, and they get the mirror. If their mirror is worth a lot, that's okay. But if they value books more, they might not get much. But since the proposals need to be complementary, if I take the bucket and all books, they can only take the mirror. So that's a valid split. But maybe the other player would prefer to have some books. Hmm. Let me think. If I take the bucket and 3 books, I get 4 +3=7, and they get 2 books and the mirror. If their mirror is worth, say, 3, and books 1 each, they get 3 +2=5. Total combined 12. If I take bucket and 2 books (6) and they take mirror and 3 books (3+3=6), total 12. Same. If I take bucket and 1 book (5) and they take mirror and 4 books (3+4=7), total 12. So regardless of how I split the books, the} \\
            }
        }
    }
    & & \\ \\

    \theutterance \stepcounter{utterance}  
    & & \multicolumn{4}{p{0.6\linewidth}}{
        \cellcolor[rgb]{0.9,0.9,0.9}{
            \makecell[{{p{\linewidth}}}]{
                \texttt{\tiny{[GM$\rangle$P2]}}
                \texttt{<think>} \\
\texttt{Okay, let's try to figure out the best way to split the items. The available items are 1 bucket, 5 books, and 1 mirror. My value function is bucket:4, book:1, mirror:1. The other player's values are secret, so I need to make a proposal that likely maximizes the total score without knowing their preferences.} \\
\\ 
\texttt{First, I should consider that the other player might value the bucket, mirror, or books more. Since there's only one bucket and one mirror, but five books, it's possible they might want the bucket or mirror. My highest value is the bucket, so I want that. But if the other player also wants the bucket, we can't both have it. So maybe I should propose taking the bucket and some books, and let them take the mirror and the rest of the books. That way, we don't conflict on the bucket and mirror. Let's say I take the bucket and 2 books (4 + 2 = 6 points), and they take the mirror and 3 books. If their mirror is worth more to them than the books, that's good. If not, maybe they'd prefer more books. But since I don't know their values, this split might be safe. Alternatively, if I take the bucket and all 5 books, I get 4 +5=9, and they get the mirror. If their mirror is worth a lot, that's okay. But if they value books more, they might not get much. But since the proposals need to be complementary, if I take the bucket and all books, they can only take the mirror. So that's a valid split. But maybe the other player would prefer to have some books. Hmm. Let me think. If I take the bucket and 3 books, I get 4 +3=7, and they get 2 books and the mirror. If their mirror is worth, say, 3, and books 1 each, they get 3 +2=5. Total combined 12. If I take bucket and 2 books (6) and they take mirror and 3 books (3+3=6), total 12. Same. If I take bucket and 1 book (5) and they take mirror and 4 books (3+4=7), total 12. So regardless of how I split the books, the} \\
            }
        }
    }
    & \\ \\

    \theutterance \stepcounter{utterance}  
    & & & \multicolumn{4}{p{0.6\linewidth}}{
        \cellcolor[rgb]{0.9,0.9,0.9}{
            \makecell[{{p{\linewidth}}}]{
                \texttt{\tiny{[GM$\langle$P2]}}
                \texttt{<think>} \\
\texttt{Okay, let's try to figure out the best way to split the items. The available items are 1 bucket, 5 books, and 1 mirror. My value function is bucket:4, book:1, mirror:1. The other player's values are secret, so I need to make a proposal that likely maximizes the total score without knowing their preferences.} \\
\\ 
\texttt{First, I should consider that the other player might value the bucket, mirror, or books more. Since there's only one bucket and one mirror, but five books, it's possible they might want the bucket or mirror. My highest value is the bucket, so I want that. But if the other player also wants the bucket, we can't both have it. So maybe I should propose taking the bucket and some books, and let them take the mirror and the rest of the books. That way, we don't conflict on the bucket and mirror. Let's say I take the bucket and 2 books (4 + 2 = 6 points), and they take the mirror and 3 books. If their mirror is worth more to them than the books, that's good. If not, maybe they'd prefer more books. But since I don't know their values, this split might be safe. Alternatively, if I take the bucket and all 5 books, I get 4 +5=9, and they get the mirror. If their mirror is worth a lot, that's okay. But if they value books more, they might not get much. But since the proposals need to be complementary, if I take the bucket and all books, they can only take the mirror. So that's a valid split. But maybe the other player would prefer to have some books. Hmm. Let me think. If I take the bucket and 3 books, I get 4 +3=7, and they get 2 books and the mirror. If their mirror is worth, say, 3, and books 1 each, they get 3 +2=5. Total combined 12. If I take bucket and 2 books (6) and they take mirror and 3 books (3+3=6), total 12. Same. If I take bucket and 1 book (5) and they take mirror and 4 books (3+4=7), total 12. So regardless of how I split the books, the} \\
            }
        }
    }
     \\ \\

    \theutterance \stepcounter{utterance}  
    & & & \multicolumn{2}{p{0.3\linewidth}}{
        \cellcolor[rgb]{0.9,0.9,0.9}{
            \makecell[{{p{\linewidth}}}]{
                \texttt{\tiny{[GM$|$GM]}}
                \texttt{continue} \\
            }
        }
    }
    & & \\ \\

    \theutterance \stepcounter{utterance}  
    & & & \multicolumn{2}{p{0.3\linewidth}}{
        \cellcolor[rgb]{0.9,0.9,0.9}{
            \makecell[{{p{\linewidth}}}]{
                \texttt{\tiny{[GM$|$GM]}}
                \texttt{<think>} \\
\texttt{Okay, let's try to figure out the best way to split the items. The available items are 1 bucket, 5 books, and 1 mirror. My value function is bucket:4, book:1, mirror:1. The other player's values are secret, so I need to make a proposal that likely maximizes the total score without knowing their preferences.} \\
\\ 
\texttt{First, I should consider that the other player might value the bucket, mirror, or books more. Since there's only one bucket and one mirror, but five books, it's possible they might want the bucket or mirror. My highest value is the bucket, so I want that. But if the other player also wants the bucket, we can't both have it. So maybe I should propose taking the bucket and some books, and let them take the mirror and the rest of the books. That way, we don't conflict on the bucket and mirror. Let's say I take the bucket and 2 books (4 + 2 = 6 points), and they take the mirror and 3 books. If their mirror is worth more to them than the books, that's good. If not, maybe they'd prefer more books. But since I don't know their values, this split might be safe. Alternatively, if I take the bucket and all 5 books, I get 4 +5=9, and they get the mirror. If their mirror is worth a lot, that's okay. But if they value books more, they might not get much. But since the proposals need to be complementary, if I take the bucket and all books, they can only take the mirror. So that's a valid split. But maybe the other player would prefer to have some books. Hmm. Let me think. If I take the bucket and 3 books, I get 4 +3=7, and they get 2 books and the mirror. If their mirror is worth, say, 3, and books 1 each, they get 3 +2=5. Total combined 12. If I take bucket and 2 books (6) and they take mirror and 3 books (3+3=6), total 12. Same. If I take bucket and 1 book (5) and they take mirror and 4 books (3+4=7), total 12. So regardless of how I split the books, the} \\
            }
        }
    }
    & & \\ \\

    \theutterance \stepcounter{utterance}  
    & & \multicolumn{4}{p{0.6\linewidth}}{
        \cellcolor[rgb]{0.9,0.9,0.9}{
            \makecell[{{p{\linewidth}}}]{
                \texttt{\tiny{[P1$\langle$GM]}}
                \texttt{<think>} \\
\texttt{Okay, let's try to figure out the best way to split the items. The available items are 1 bucket, 5 books, and 1 mirror. My value function is bucket:4, book:1, mirror:1. The other player's values are secret, so I need to make a proposal that likely maximizes the total score without knowing their preferences.} \\
\\ 
\texttt{First, I should consider that the other player might value the bucket, mirror, or books more. Since there's only one bucket and one mirror, but five books, it's possible they might want the bucket or mirror. My highest value is the bucket, so I want that. But if the other player also wants the bucket, we can't both have it. So maybe I should propose taking the bucket and some books, and let them take the mirror and the rest of the books. That way, we don't conflict on the bucket and mirror. Let's say I take the bucket and 2 books (4 + 2 = 6 points), and they take the mirror and 3 books. If their mirror is worth more to them than the books, that's good. If not, maybe they'd prefer more books. But since I don't know their values, this split might be safe. Alternatively, if I take the bucket and all 5 books, I get 4 +5=9, and they get the mirror. If their mirror is worth a lot, that's okay. But if they value books more, they might not get much. But since the proposals need to be complementary, if I take the bucket and all books, they can only take the mirror. So that's a valid split. But maybe the other player would prefer to have some books. Hmm. Let me think. If I take the bucket and 3 books, I get 4 +3=7, and they get 2 books and the mirror. If their mirror is worth, say, 3, and books 1 each, they get 3 +2=5. Total combined 12. If I take bucket and 2 books (6) and they take mirror and 3 books (3+3=6), total 12. Same. If I take bucket and 1 book (5) and they take mirror and 4 books (3+4=7), total 12. So regardless of how I split the books, the} \\
            }
        }
    }
    & \\ \\

    \theutterance \stepcounter{utterance}  
    & \multicolumn{4}{p{0.6\linewidth}}{
        \cellcolor[rgb]{0.9,0.9,0.9}{
            \makecell[{{p{\linewidth}}}]{
                \texttt{\tiny{[P1$\rangle$GM]}}
                \texttt{<think>} \\
\texttt{Okay, let's try to figure out the best way to split the items. The available items are 1 bucket, 5 books, and 1 mirror. My value function is bucket:4, book:1, mirror:1. The other player's values are secret, so I need to make a proposal that likely maximizes the total score without knowing their preferences.} \\
\\ 
\texttt{First, I should consider that the other player might value the bucket, mirror, or books more. Since there's only one bucket and one mirror, but five books, it's possible they might want the bucket or mirror. My highest value is the bucket, so I want that. But if the other player also wants the bucket, we can't both have it. So maybe I should propose taking the bucket and some books, and let them take the mirror and the rest of the books. That way, we don't conflict on the bucket and mirror. Let's say I take the bucket and 2 books (4 + 2 = 6 points), and they take the mirror and 3 books. If their mirror is worth more to them than the books, that's good. If not, maybe they'd prefer more books. But since I don't know their values, this split might be safe. Alternatively, if I take the bucket and all 5 books, I get 4 +5=9, and they get the mirror. If their mirror is worth a lot, that's okay. But if they value books more, they might not get much. But since the proposals need to be complementary, if I take the bucket and all books, they can only take the mirror. So that's a valid split. But maybe the other player would prefer to have some books. Hmm. Let me think. If I take the bucket and 3 books, I get 4 +3=7, and they get 2 books and the mirror. If their mirror is worth, say, 3, and books 1 each, they get 3 +2=5. Total combined 12. If I take bucket and 2 books (6) and they take mirror and 3 books (3+3=6), total 12. Same. If I take bucket and 1 book (5) and they take mirror and 4 books (3+4=7), total 12. So regardless of how I split the books, the} \\
            }
        }
    }
    & & \\ \\

    \theutterance \stepcounter{utterance}  
    & & & \multicolumn{2}{p{0.3\linewidth}}{
        \cellcolor[rgb]{0.9,0.9,0.9}{
            \makecell[{{p{\linewidth}}}]{
                \texttt{\tiny{[GM$|$GM]}}
                \texttt{continue} \\
            }
        }
    }
    & & \\ \\

    \theutterance \stepcounter{utterance}  
    & & & \multicolumn{2}{p{0.3\linewidth}}{
        \cellcolor[rgb]{0.9,0.9,0.9}{
            \makecell[{{p{\linewidth}}}]{
                \texttt{\tiny{[GM$|$GM]}}
                \texttt{<think>} \\
\texttt{Okay, let's try to figure out the best way to split the items. The available items are 1 bucket, 5 books, and 1 mirror. My value function is bucket:4, book:1, mirror:1. The other player's values are secret, so I need to make a proposal that likely maximizes the total score without knowing their preferences.} \\
\\ 
\texttt{First, I should consider that the other player might value the bucket, mirror, or books more. Since there's only one bucket and one mirror, but five books, it's possible they might want the bucket or mirror. My highest value is the bucket, so I want that. But if the other player also wants the bucket, we can't both have it. So maybe I should propose taking the bucket and some books, and let them take the mirror and the rest of the books. That way, we don't conflict on the bucket and mirror. Let's say I take the bucket and 2 books (4 + 2 = 6 points), and they take the mirror and 3 books. If their mirror is worth more to them than the books, that's good. If not, maybe they'd prefer more books. But since I don't know their values, this split might be safe. Alternatively, if I take the bucket and all 5 books, I get 4 +5=9, and they get the mirror. If their mirror is worth a lot, that's okay. But if they value books more, they might not get much. But since the proposals need to be complementary, if I take the bucket and all books, they can only take the mirror. So that's a valid split. But maybe the other player would prefer to have some books. Hmm. Let me think. If I take the bucket and 3 books, I get 4 +3=7, and they get 2 books and the mirror. If their mirror is worth, say, 3, and books 1 each, they get 3 +2=5. Total combined 12. If I take bucket and 2 books (6) and they take mirror and 3 books (3+3=6), total 12. Same. If I take bucket and 1 book (5) and they take mirror and 4 books (3+4=7), total 12. So regardless of how I split the books, the} \\
            }
        }
    }
    & & \\ \\

    \theutterance \stepcounter{utterance}  
    & & \multicolumn{4}{p{0.6\linewidth}}{
        \cellcolor[rgb]{0.9,0.9,0.9}{
            \makecell[{{p{\linewidth}}}]{
                \texttt{\tiny{[GM$\rangle$P2]}}
                \texttt{<think>} \\
\texttt{Okay, let's try to figure out the best way to split the items. The available items are 1 bucket, 5 books, and 1 mirror. My value function is bucket:4, book:1, mirror:1. The other player's values are secret, so I need to make a proposal that likely maximizes the total score without knowing their preferences.} \\
\\ 
\texttt{First, I should consider that the other player might value the bucket, mirror, or books more. Since there's only one bucket and one mirror, but five books, it's possible they might want the bucket or mirror. My highest value is the bucket, so I want that. But if the other player also wants the bucket, we can't both have it. So maybe I should propose taking the bucket and some books, and let them take the mirror and the rest of the books. That way, we don't conflict on the bucket and mirror. Let's say I take the bucket and 2 books (4 + 2 = 6 points), and they take the mirror and 3 books. If their mirror is worth more to them than the books, that's good. If not, maybe they'd prefer more books. But since I don't know their values, this split might be safe. Alternatively, if I take the bucket and all 5 books, I get 4 +5=9, and they get the mirror. If their mirror is worth a lot, that's okay. But if they value books more, they might not get much. But since the proposals need to be complementary, if I take the bucket and all books, they can only take the mirror. So that's a valid split. But maybe the other player would prefer to have some books. Hmm. Let me think. If I take the bucket and 3 books, I get 4 +3=7, and they get 2 books and the mirror. If their mirror is worth, say, 3, and books 1 each, they get 3 +2=5. Total combined 12. If I take bucket and 2 books (6) and they take mirror and 3 books (3+3=6), total 12. Same. If I take bucket and 1 book (5) and they take mirror and 4 books (3+4=7), total 12. So regardless of how I split the books, the} \\
            }
        }
    }
    & \\ \\

    \theutterance \stepcounter{utterance}  
    & & & \multicolumn{4}{p{0.6\linewidth}}{
        \cellcolor[rgb]{0.9,0.9,0.9}{
            \makecell[{{p{\linewidth}}}]{
                \texttt{\tiny{[GM$\langle$P2]}}
                \texttt{<think>} \\
\texttt{Okay, let's try to figure out the best way to split the items. The available items are 1 bucket, 5 books, and 1 mirror. My value function is bucket:4, book:1, mirror:1. The other player's values are secret, so I need to make a proposal that likely maximizes the total score without knowing their preferences.} \\
\\ 
\texttt{First, I should consider that the other player might value the bucket, mirror, or books more. Since there's only one bucket and one mirror, but five books, it's possible they might want the bucket or mirror. My highest value is the bucket, so I want that. But if the other player also wants the bucket, we can't both have it. So maybe I should propose taking the bucket and some books, and let them take the mirror and the rest of the books. That way, we don't conflict on the bucket and mirror. Let's say I take the bucket and 2 books (4 + 2 = 6 points), and they take the mirror and 3 books. If their mirror is worth more to them than the books, that's good. If not, maybe they'd prefer more books. But since I don't know their values, this split might be safe. Alternatively, if I take the bucket and all 5 books, I get 4 +5=9, and they get the mirror. If their mirror is worth a lot, that's okay. But if they value books more, they might not get much. But since the proposals need to be complementary, if I take the bucket and all books, they can only take the mirror. So that's a valid split. But maybe the other player would prefer to have some books. Hmm. Let me think. If I take the bucket and 3 books, I get 4 +3=7, and they get 2 books and the mirror. If their mirror is worth, say, 3, and books 1 each, they get 3 +2=5. Total combined 12. If I take bucket and 2 books (6) and they take mirror and 3 books (3+3=6), total 12. Same. If I take bucket and 1 book (5) and they take mirror and 4 books (3+4=7), total 12. So regardless of how I split the books, the} \\
            }
        }
    }
     \\ \\

    \theutterance \stepcounter{utterance}  
    & & & \multicolumn{2}{p{0.3\linewidth}}{
        \cellcolor[rgb]{0.9,0.9,0.9}{
            \makecell[{{p{\linewidth}}}]{
                \texttt{\tiny{[GM$|$GM]}}
                \texttt{continue} \\
            }
        }
    }
    & & \\ \\

    \theutterance \stepcounter{utterance}  
    & & & \multicolumn{2}{p{0.3\linewidth}}{
        \cellcolor[rgb]{0.9,0.9,0.9}{
            \makecell[{{p{\linewidth}}}]{
                \texttt{\tiny{[GM$|$GM]}}
                \texttt{<think>} \\
\texttt{Okay, let's try to figure out the best way to split the items. The available items are 1 bucket, 5 books, and 1 mirror. My value function is bucket:4, book:1, mirror:1. The other player's values are secret, so I need to make a proposal that likely maximizes the total score without knowing their preferences.} \\
\\ 
\texttt{First, I should consider that the other player might value the bucket, mirror, or books more. Since there's only one bucket and one mirror, but five books, it's possible they might want the bucket or mirror. My highest value is the bucket, so I want that. But if the other player also wants the bucket, we can't both have it. So maybe I should propose taking the bucket and some books, and let them take the mirror and the rest of the books. That way, we don't conflict on the bucket and mirror. Let's say I take the bucket and 2 books (4 + 2 = 6 points), and they take the mirror and 3 books. If their mirror is worth more to them than the books, that's good. If not, maybe they'd prefer more books. But since I don't know their values, this split might be safe. Alternatively, if I take the bucket and all 5 books, I get 4 +5=9, and they get the mirror. If their mirror is worth a lot, that's okay. But if they value books more, they might not get much. But since the proposals need to be complementary, if I take the bucket and all books, they can only take the mirror. So that's a valid split. But maybe the other player would prefer to have some books. Hmm. Let me think. If I take the bucket and 3 books, I get 4 +3=7, and they get 2 books and the mirror. If their mirror is worth, say, 3, and books 1 each, they get 3 +2=5. Total combined 12. If I take bucket and 2 books (6) and they take mirror and 3 books (3+3=6), total 12. Same. If I take bucket and 1 book (5) and they take mirror and 4 books (3+4=7), total 12. So regardless of how I split the books, the} \\
            }
        }
    }
    & & \\ \\

    \theutterance \stepcounter{utterance}  
    & & \multicolumn{4}{p{0.6\linewidth}}{
        \cellcolor[rgb]{0.9,0.9,0.9}{
            \makecell[{{p{\linewidth}}}]{
                \texttt{\tiny{[P1$\langle$GM]}}
                \texttt{<think>} \\
\texttt{Okay, let's try to figure out the best way to split the items. The available items are 1 bucket, 5 books, and 1 mirror. My value function is bucket:4, book:1, mirror:1. The other player's values are secret, so I need to make a proposal that likely maximizes the total score without knowing their preferences.} \\
\\ 
\texttt{First, I should consider that the other player might value the bucket, mirror, or books more. Since there's only one bucket and one mirror, but five books, it's possible they might want the bucket or mirror. My highest value is the bucket, so I want that. But if the other player also wants the bucket, we can't both have it. So maybe I should propose taking the bucket and some books, and let them take the mirror and the rest of the books. That way, we don't conflict on the bucket and mirror. Let's say I take the bucket and 2 books (4 + 2 = 6 points), and they take the mirror and 3 books. If their mirror is worth more to them than the books, that's good. If not, maybe they'd prefer more books. But since I don't know their values, this split might be safe. Alternatively, if I take the bucket and all 5 books, I get 4 +5=9, and they get the mirror. If their mirror is worth a lot, that's okay. But if they value books more, they might not get much. But since the proposals need to be complementary, if I take the bucket and all books, they can only take the mirror. So that's a valid split. But maybe the other player would prefer to have some books. Hmm. Let me think. If I take the bucket and 3 books, I get 4 +3=7, and they get 2 books and the mirror. If their mirror is worth, say, 3, and books 1 each, they get 3 +2=5. Total combined 12. If I take bucket and 2 books (6) and they take mirror and 3 books (3+3=6), total 12. Same. If I take bucket and 1 book (5) and they take mirror and 4 books (3+4=7), total 12. So regardless of how I split the books, the} \\
            }
        }
    }
    & \\ \\

    \theutterance \stepcounter{utterance}  
    & \multicolumn{4}{p{0.6\linewidth}}{
        \cellcolor[rgb]{0.9,0.9,0.9}{
            \makecell[{{p{\linewidth}}}]{
                \texttt{\tiny{[P1$\rangle$GM]}}
                \texttt{<think>} \\
\texttt{Okay, let's try to figure out the best way to split the items. The available items are 1 bucket, 5 books, and 1 mirror. My value function is bucket:4, book:1, mirror:1. The other player's values are secret, so I need to make a proposal that likely maximizes the total score without knowing their preferences.} \\
\\ 
\texttt{First, I should consider that the other player might value the bucket, mirror, or books more. Since there's only one bucket and one mirror, but five books, it's possible they might want the bucket or mirror. My highest value is the bucket, so I want that. But if the other player also wants the bucket, we can't both have it. So maybe I should propose taking the bucket and some books, and let them take the mirror and the rest of the books. That way, we don't conflict on the bucket and mirror. Let's say I take the bucket and 2 books (4 + 2 = 6 points), and they take the mirror and 3 books. If their mirror is worth more to them than the books, that's good. If not, maybe they'd prefer more books. But since I don't know their values, this split might be safe. Alternatively, if I take the bucket and all 5 books, I get 4 +5=9, and they get the mirror. If their mirror is worth a lot, that's okay. But if they value books more, they might not get much. But since the proposals need to be complementary, if I take the bucket and all books, they can only take the mirror. So that's a valid split. But maybe the other player would prefer to have some books. Hmm. Let me think. If I take the bucket and 3 books, I get 4 +3=7, and they get 2 books and the mirror. If their mirror is worth, say, 3, and books 1 each, they get 3 +2=5. Total combined 12. If I take bucket and 2 books (6) and they take mirror and 3 books (3+3=6), total 12. Same. If I take bucket and 1 book (5) and they take mirror and 4 books (3+4=7), total 12. So regardless of how I split the books, the} \\
            }
        }
    }
    & & \\ \\

    \theutterance \stepcounter{utterance}  
    & & & \multicolumn{2}{p{0.3\linewidth}}{
        \cellcolor[rgb]{0.9,0.9,0.9}{
            \makecell[{{p{\linewidth}}}]{
                \texttt{\tiny{[GM$|$GM]}}
                \texttt{continue} \\
            }
        }
    }
    & & \\ \\

    \theutterance \stepcounter{utterance}  
    & & & \multicolumn{2}{p{0.3\linewidth}}{
        \cellcolor[rgb]{0.9,0.9,0.9}{
            \makecell[{{p{\linewidth}}}]{
                \texttt{\tiny{[GM$|$GM]}}
                \texttt{<think>} \\
\texttt{Okay, let's try to figure out the best way to split the items. The available items are 1 bucket, 5 books, and 1 mirror. My value function is bucket:4, book:1, mirror:1. The other player's values are secret, so I need to make a proposal that likely maximizes the total score without knowing their preferences.} \\
\\ 
\texttt{First, I should consider that the other player might value the bucket, mirror, or books more. Since there's only one bucket and one mirror, but five books, it's possible they might want the bucket or mirror. My highest value is the bucket, so I want that. But if the other player also wants the bucket, we can't both have it. So maybe I should propose taking the bucket and some books, and let them take the mirror and the rest of the books. That way, we don't conflict on the bucket and mirror. Let's say I take the bucket and 2 books (4 + 2 = 6 points), and they take the mirror and 3 books. If their mirror is worth more to them than the books, that's good. If not, maybe they'd prefer more books. But since I don't know their values, this split might be safe. Alternatively, if I take the bucket and all 5 books, I get 4 +5=9, and they get the mirror. If their mirror is worth a lot, that's okay. But if they value books more, they might not get much. But since the proposals need to be complementary, if I take the bucket and all books, they can only take the mirror. So that's a valid split. But maybe the other player would prefer to have some books. Hmm. Let me think. If I take the bucket and 3 books, I get 4 +3=7, and they get 2 books and the mirror. If their mirror is worth, say, 3, and books 1 each, they get 3 +2=5. Total combined 12. If I take bucket and 2 books (6) and they take mirror and 3 books (3+3=6), total 12. Same. If I take bucket and 1 book (5) and they take mirror and 4 books (3+4=7), total 12. So regardless of how I split the books, the} \\
            }
        }
    }
    & & \\ \\

    \theutterance \stepcounter{utterance}  
    & & \multicolumn{4}{p{0.6\linewidth}}{
        \cellcolor[rgb]{0.9,0.9,0.9}{
            \makecell[{{p{\linewidth}}}]{
                \texttt{\tiny{[GM$\rangle$P2]}}
                \texttt{<think>} \\
\texttt{Okay, let's try to figure out the best way to split the items. The available items are 1 bucket, 5 books, and 1 mirror. My value function is bucket:4, book:1, mirror:1. The other player's values are secret, so I need to make a proposal that likely maximizes the total score without knowing their preferences.} \\
\\ 
\texttt{First, I should consider that the other player might value the bucket, mirror, or books more. Since there's only one bucket and one mirror, but five books, it's possible they might want the bucket or mirror. My highest value is the bucket, so I want that. But if the other player also wants the bucket, we can't both have it. So maybe I should propose taking the bucket and some books, and let them take the mirror and the rest of the books. That way, we don't conflict on the bucket and mirror. Let's say I take the bucket and 2 books (4 + 2 = 6 points), and they take the mirror and 3 books. If their mirror is worth more to them than the books, that's good. If not, maybe they'd prefer more books. But since I don't know their values, this split might be safe. Alternatively, if I take the bucket and all 5 books, I get 4 +5=9, and they get the mirror. If their mirror is worth a lot, that's okay. But if they value books more, they might not get much. But since the proposals need to be complementary, if I take the bucket and all books, they can only take the mirror. So that's a valid split. But maybe the other player would prefer to have some books. Hmm. Let me think. If I take the bucket and 3 books, I get 4 +3=7, and they get 2 books and the mirror. If their mirror is worth, say, 3, and books 1 each, they get 3 +2=5. Total combined 12. If I take bucket and 2 books (6) and they take mirror and 3 books (3+3=6), total 12. Same. If I take bucket and 1 book (5) and they take mirror and 4 books (3+4=7), total 12. So regardless of how I split the books, the} \\
            }
        }
    }
    & \\ \\

    \theutterance \stepcounter{utterance}  
    & & & \multicolumn{4}{p{0.6\linewidth}}{
        \cellcolor[rgb]{0.9,0.9,0.9}{
            \makecell[{{p{\linewidth}}}]{
                \texttt{\tiny{[GM$\langle$P2]}}
                \texttt{<think>} \\
\texttt{Okay, let's try to figure out the best way to split the items. The available items are 1 bucket, 5 books, and 1 mirror. My value function is bucket:4, book:1, mirror:1. The other player's values are secret, so I need to make a proposal that likely maximizes the total score without knowing their preferences.} \\
\\ 
\texttt{First, I should consider that the other player might value the bucket, mirror, or books more. Since there's only one bucket and one mirror, but five books, it's possible they might want the bucket or mirror. My highest value is the bucket, so I want that. But if the other player also wants the bucket, we can't both have it. So maybe I should propose taking the bucket and some books, and let them take the mirror and the rest of the books. That way, we don't conflict on the bucket and mirror. Let's say I take the bucket and 2 books (4 + 2 = 6 points), and they take the mirror and 3 books. If their mirror is worth more to them than the books, that's good. If not, maybe they'd prefer more books. But since I don't know their values, this split might be safe. Alternatively, if I take the bucket and all 5 books, I get 4 +5=9, and they get the mirror. If their mirror is worth a lot, that's okay. But if they value books more, they might not get much. But since the proposals need to be complementary, if I take the bucket and all books, they can only take the mirror. So that's a valid split. But maybe the other player would prefer to have some books. Hmm. Let me think. If I take the bucket and 3 books, I get 4 +3=7, and they get 2 books and the mirror. If their mirror is worth, say, 3, and books 1 each, they get 3 +2=5. Total combined 12. If I take bucket and 2 books (6) and they take mirror and 3 books (3+3=6), total 12. Same. If I take bucket and 1 book (5) and they take mirror and 4 books (3+4=7), total 12. So regardless of how I split the books, the} \\
            }
        }
    }
     \\ \\

    \theutterance \stepcounter{utterance}  
    & & & \multicolumn{2}{p{0.3\linewidth}}{
        \cellcolor[rgb]{0.9,0.9,0.9}{
            \makecell[{{p{\linewidth}}}]{
                \texttt{\tiny{[GM$|$GM]}}
                \texttt{continue} \\
            }
        }
    }
    & & \\ \\

    \theutterance \stepcounter{utterance}  
    & & & \multicolumn{2}{p{0.3\linewidth}}{
        \cellcolor[rgb]{0.9,0.9,0.9}{
            \makecell[{{p{\linewidth}}}]{
                \texttt{\tiny{[GM$|$GM]}}
                \texttt{<think>} \\
\texttt{Okay, let's try to figure out the best way to split the items. The available items are 1 bucket, 5 books, and 1 mirror. My value function is bucket:4, book:1, mirror:1. The other player's values are secret, so I need to make a proposal that likely maximizes the total score without knowing their preferences.} \\
\\ 
\texttt{First, I should consider that the other player might value the bucket, mirror, or books more. Since there's only one bucket and one mirror, but five books, it's possible they might want the bucket or mirror. My highest value is the bucket, so I want that. But if the other player also wants the bucket, we can't both have it. So maybe I should propose taking the bucket and some books, and let them take the mirror and the rest of the books. That way, we don't conflict on the bucket and mirror. Let's say I take the bucket and 2 books (4 + 2 = 6 points), and they take the mirror and 3 books. If their mirror is worth more to them than the books, that's good. If not, maybe they'd prefer more books. But since I don't know their values, this split might be safe. Alternatively, if I take the bucket and all 5 books, I get 4 +5=9, and they get the mirror. If their mirror is worth a lot, that's okay. But if they value books more, they might not get much. But since the proposals need to be complementary, if I take the bucket and all books, they can only take the mirror. So that's a valid split. But maybe the other player would prefer to have some books. Hmm. Let me think. If I take the bucket and 3 books, I get 4 +3=7, and they get 2 books and the mirror. If their mirror is worth, say, 3, and books 1 each, they get 3 +2=5. Total combined 12. If I take bucket and 2 books (6) and they take mirror and 3 books (3+3=6), total 12. Same. If I take bucket and 1 book (5) and they take mirror and 4 books (3+4=7), total 12. So regardless of how I split the books, the} \\
            }
        }
    }
    & & \\ \\

    \theutterance \stepcounter{utterance}  
    & & \multicolumn{4}{p{0.6\linewidth}}{
        \cellcolor[rgb]{0.9,0.9,0.9}{
            \makecell[{{p{\linewidth}}}]{
                \texttt{\tiny{[P1$\langle$GM]}}
                \texttt{<think>} \\
\texttt{Okay, let's try to figure out the best way to split the items. The available items are 1 bucket, 5 books, and 1 mirror. My value function is bucket:4, book:1, mirror:1. The other player's values are secret, so I need to make a proposal that likely maximizes the total score without knowing their preferences.} \\
\\ 
\texttt{First, I should consider that the other player might value the bucket, mirror, or books more. Since there's only one bucket and one mirror, but five books, it's possible they might want the bucket or mirror. My highest value is the bucket, so I want that. But if the other player also wants the bucket, we can't both have it. So maybe I should propose taking the bucket and some books, and let them take the mirror and the rest of the books. That way, we don't conflict on the bucket and mirror. Let's say I take the bucket and 2 books (4 + 2 = 6 points), and they take the mirror and 3 books. If their mirror is worth more to them than the books, that's good. If not, maybe they'd prefer more books. But since I don't know their values, this split might be safe. Alternatively, if I take the bucket and all 5 books, I get 4 +5=9, and they get the mirror. If their mirror is worth a lot, that's okay. But if they value books more, they might not get much. But since the proposals need to be complementary, if I take the bucket and all books, they can only take the mirror. So that's a valid split. But maybe the other player would prefer to have some books. Hmm. Let me think. If I take the bucket and 3 books, I get 4 +3=7, and they get 2 books and the mirror. If their mirror is worth, say, 3, and books 1 each, they get 3 +2=5. Total combined 12. If I take bucket and 2 books (6) and they take mirror and 3 books (3+3=6), total 12. Same. If I take bucket and 1 book (5) and they take mirror and 4 books (3+4=7), total 12. So regardless of how I split the books, the} \\
\\ 
\\ 
\texttt{The time is up. It is now your turn to submit a secret proposal.} \\
            }
        }
    }
    & \\ \\

    \theutterance \stepcounter{utterance}  
    & \multicolumn{4}{p{0.6\linewidth}}{
        \cellcolor[rgb]{0.9,0.9,0.9}{
            \makecell[{{p{\linewidth}}}]{
                \texttt{\tiny{[P1$\rangle$GM]}}
                \texttt{<think>} \\
\texttt{Okay, let's try to figure out the best way to split the items. The available items are 1 bucket, 5 books, and 1 mirror. My value function is bucket:4, book:1, mirror:1. The other player's values are secret, so I need to make a proposal that likely maximizes the total score without knowing their preferences.} \\
\\ 
\texttt{First, I should consider that the other player might value the bucket, mirror, or books more. Since there's only one bucket and one mirror, but five books, it's possible they might want the bucket or mirror. My highest value is the bucket, so I want that. But if the other player also wants the bucket, we can't both have it. So maybe I should propose taking the bucket and some books, and let them take the mirror and the rest of the books. That way, we don't conflict on the bucket and mirror. Let's say I take the bucket and 2 books (4 + 2 = 6 points), and they take the mirror and 3 books. If their mirror is worth more to them than the books, that's good. If not, maybe they'd prefer more books. But since I don't know their values, this split might be safe. Alternatively, if I take the bucket and all 5 books, I get 4 +5=9, and they get the mirror. If their mirror is worth a lot, that's okay. But if they value books more, they might not get much. But since the proposals need to be complementary, if I take the bucket and all books, they can only take the mirror. So that's a valid split. But maybe the other player would prefer to have some books. Hmm. Let me think. If I take the bucket and 3 books, I get 4 +3=7, and they get 2 books and the mirror. If their mirror is worth, say, 3, and books 1 each, they get 3 +2=5. Total combined 12. If I take bucket and 2 books (6) and they take mirror and 3 books (3+3=6), total 12. Same. If I take bucket and 1 book (5) and they take mirror and 4 books (3+4=7), total 12. So regardless of how I split the books, the} \\
            }
        }
    }
    & & \\ \\

    \theutterance \stepcounter{utterance}  
    & & & \multicolumn{2}{p{0.3\linewidth}}{
        \cellcolor[rgb]{0.9,0.9,0.9}{
            \makecell[{{p{\linewidth}}}]{
                \texttt{\tiny{[GM$|$GM]}}
                \texttt{continue} \\
            }
        }
    }
    & & \\ \\

    \theutterance \stepcounter{utterance}  
    & & & \multicolumn{2}{p{0.3\linewidth}}{
        \cellcolor[rgb]{0.9,0.9,0.9}{
            \makecell[{{p{\linewidth}}}]{
                \texttt{\tiny{[GM$|$GM]}}
                \texttt{player was instructed to make a proposal} \\
            }
        }
    }
    & & \\ \\

\end{supertabular}
}

\end{document}
